% Table 4: Diagnostic Categorization of Clean-Background Test Alarms
\begin{table}[H]
\caption{Rule-based diagnostic categorization of a stratified random sample of $N = 100$ clean-background test alarms raised across the five pilot stations (2023--2025), disaggregated across detector families. Sub-sample counts reflect multi-model alarm overlaps ($N=45$ for gross count baselines, $N=41$ for Tier 3 Weather Fusion, and $N=37$ for Tier 4 Physics-Motivated). Confidence intervals are 95\% Wilson score intervals for the overall sample. Note that active precipitation accounts for 15.6\% of baseline alarms and 14.6\% of Tier 3 alarms, but only 2.7\% (1 alarm) for Tier 4.}
\label{tab:alarm_adjudication}
\centering
\small
\begin{tabularx}{\textwidth}{@{}X c c c c c@{}}
\toprule
\textbf{Adjudicated Origin} & \textbf{Baselines ($N=45$)} & \textbf{Tier 3 ($N=41$)} & \textbf{Tier 4 ($N=37$)} & \textbf{Overall ($N=100$)} & \textbf{95\% Wilson CI} \\
\midrule
Active Precipitation Radon Washout & 7 (15.6\%) & 6 (14.6\%) & 1 (2.7\%) & 12 (12.0\%) & [7.0\%, 19.8\%] \\
Post-Precipitation Progeny Residue ($<$ 6h) & 0 (0.0\%) & 1 (2.4\%) & 0 (0.0\%) & 2 (2.0\%) & [0.6\%, 7.0\%] \\
Particulate Filter Advance / Step Artifact & 2 (4.4\%) & 0 (0.0\%) & 1 (2.7\%) & 2 (2.0\%) & [0.6\%, 7.0\%] \\
Instrument Baseline Step / Gain Drift (Tampa 2023) & 8 (17.8\%) & 18 (43.9\%) & 12 (32.4\%) & 28 (28.0\%) & [20.1\%, 37.5\%] \\
Other Documented Natural Phenomena & 0 (0.0\%) & 0 (0.0\%) & 0 (0.0\%) & 0 (0.0\%) & [0.0\%, 3.7\%] \\
Unexplained Statistical Background Fluctuations & 28 (62.2\%) & 16 (39.0\%) & 23 (62.2\%) & 56 (56.0\%) & [46.2\%, 65.3\%] \\
\midrule
\textbf{Physical / Technical Cause} & \textbf{17 (37.8\%)} & \textbf{25 (61.0\%)} & \textbf{14 (37.8\%)} & \textbf{44 (44.0\%)} & \textbf{[34.7\%, 53.8\%]} \\
\textbf{Pure Statistical Noise} & \textbf{28 (62.2\%)} & \textbf{16 (39.0\%)} & \textbf{23 (62.2\%)} & \textbf{56 (56.0\%)} & \textbf{[46.2\%, 65.3\%]} \\
\bottomrule
\end{tabularx}
\end{table}
